\begin{codebox}
\codeline{\textcolor{cmtgray}{\#\ ==============================}}
\codeline{\textcolor{cmtgray}{\#\ 3.\ 两者的区别}}
\codeline{\textcolor{cmtgray}{\#\ ==============================}}
\codeline{\textcolor{cmtgray}{\#\ \(\exists\)R.C（存在限制）：至少一个，即"存在"}}
\codeline{\textcolor{cmtgray}{\#\ \ \ \(\rightarrow\)\ 如果满足，推理器保证存在一个相应个体}}
\codeline{\textcolor{cmtgray}{\#\ \ \ \(\rightarrow\)\ 适合表达"必须有..."}}
\codeline{\textcolor{cmtgray}{\#}}
\codeline{\textcolor{cmtgray}{\#\ \(\forall\)R.C（全称限制）：全部满足，即"所有"}}
\codeline{\textcolor{cmtgray}{\#\ \ \ \(\rightarrow\)\ 如果有关系存在，则所有关系对象都必须满足约束}}
\codeline{\textcolor{cmtgray}{\#\ \ \ \(\rightarrow\)\ 如果没有关系，则全称限制自动满足（空满足）}}
\codeline{\textcolor{cmtgray}{\#\ \ \ \(\rightarrow\)\ 适合表达"只允许..."}}
\codeblank{}
\end{codebox}
